% MATERI INSTRUKSIONAL INDEPENDEN --- BUKAN TEKS SUMBER WEN-WEI LI.
% stable-unit-id: o014.mastery.diagram-chasing.001
% provenance: OpenAI Codex gpt-5.6-sol, Ultra
% license: CC BY 4.0; kompatibel dengan lisensi terjemahan sumber.

\section*{Jembatan Penguasaan 001: Pengejaran Diagram}
\addcontentsline{toc}{section}{Jembatan Penguasaan 001: Pengejaran Diagram}
\hypertarget{o014.mastery.diagram-chasing.001}{ }
\label{mastery:diagram-chasing-001}

\noindent\textbf{Status dan atribusi.}
Unit ini merupakan materi instruksional independen yang disusun untuk
mendampingi buku ini. Unit ini bukan bagian dari teks sumber Wen-Wei Li dan
tidak boleh dinisbatkan kepadanya. Provenans penyusunan:
\texttt{OpenAI Codex gpt-5.6-sol, Ultra}. Materi ini tersedia berdasarkan
Creative Commons Attribution 4.0 International
(CC BY 4.0; \url{https://creativecommons.org/licenses/by/4.0/}), sehingga
kompatibel dengan lisensi teks sumber.

\subsection*{Peta prasyarat}

Unit ini dapat dibaca setelah menguasai rangkaian prasyarat singkat berikut.
\begin{compactitem}
	\item Kernel, kokernel, dan funktorialitasnya:
	\S\ref{sec:Ker-Coker}, khususnya
	\eqref{eqn:Ker-Coker-functoriality}.
	\item Citra, kocitra, dan faktorisasi epi--mono:
	\S\ref{sec:images}.
	\item Kategori abelian serta kestabilan epimorfisme dan monomorfisme
	terhadap tarik balik dan dorong keluar:
	\S\ref{sec:Abel-cat-def} dan
	Proposisi~\ref{prop:Abel-cat-pull-push}.
	\item Kompleks, kohomologi, dan keeksakan:
	\S\ref{sec:cohomology}, terutama
	Definisi~\ref{def:cohomology}.
	\item Konstruksi kategoris morfisme penghubung dan Lema Ular:
	\S\ref{sec:diagram-lemmas}, Catatan~\ref{rem:connecting-canonical}, dan
	Teorema~\ref{prop:snake-lemma}.
\end{compactitem}
Secara konseptual, alurnya ialah
\[
	\text{kernel/kokernel}
	\Longrightarrow \text{citra dan keeksakan}
	\Longrightarrow \text{angkat--dorong--turunkan}
	\Longrightarrow \text{morfisme penghubung}.
\]

\subsection*{Metode kerja: mulai dari hambatan}

Untuk pengejaran unsur dalam kategori modul, gunakan disiplin berikut.
\begin{enumerate}
	\item Tuliskan tepat apa yang harus dibuktikan: keanggotaan dalam kernel,
	citra, atau kesamaan dua kelas kokernel.
	\item Mulailah dari unsur yang mengukur hambatan, biasanya unsur kernel
	atau kelas kokernel.
	\item Angkat unsur hanya melalui panah yang diketahui surjektif. Dorong
	unsur sepanjang panah mana pun yang tersedia.
	\item Gunakan komutativitas untuk memindahkan suatu persamaan ke sisi lain
	diagram, lalu gunakan keeksakan untuk mengganti pernyataan ``bernilai nol''
	dengan ``berasal dari citra sebelumnya''.
	\item Jika suatu pilihan diperlukan, buktikan bahwa selisih dua pilihan
	lenyap dalam target akhir.
\end{enumerate}
Dalam kategori abelian umum, ``unsur'' dapat diganti oleh morfisme dari objek
uji. Bila pengangkatan tidak tersedia secara langsung, gunakan penutup
epimorfik sebagaimana dalam Lema~\ref{prop:Abel-cat-exact-aux}. Dengan cara
ini, pengejaran unsur berikut bukan sekadar heuristik: pengejaran itu merupakan model
konkret bagi argumen kernel, kokernel, dan faktorisasi.

Sebagai contoh kerja, perhatikan diagram komutatif modul-$R$
\[
\begin{tikzcd}
	0 \arrow[r] & A' \arrow[r, "i"] \arrow[d, "a'"'] & A \arrow[r, "p"] \arrow[d, "a"'] & A'' \arrow[r] \arrow[d, "a''"] & 0 \\
	0 \arrow[r] & B' \arrow[r, "j"'] & B \arrow[r, "q"'] & B'' \arrow[r] & 0
\end{tikzcd}
\]
dengan kedua baris eksak. Andaikan $a'$ dan $a''$ injektif. Untuk membuktikan
$a$ injektif, mulai dengan hambatan $x\in\Ker(a)$. Komutativitas memberikan
$a''p(x)=qa(x)=0$, sehingga injektivitas $a''$ memberi $p(x)=0$. Keeksakan
baris atas memberi $x=i(x')$ untuk suatu $x'\in A'$. Selanjutnya,
\[
	0=a(x)=a i(x')=j a'(x').
\]
Karena $j$ dan $a'$ injektif, $x'=0$, jadi $x=0$. Perhatikan pola yang
berulang: mulai dari kernel, dorong ke kanan, gunakan injektivitas, kembali
ke kiri melalui keeksakan, lalu gunakan injektivitas sekali lagi.

\subsection*{Latihan dan solusi lengkap}

% stable-exercise-id: o014.mastery.diagram-chasing.001.ex01
\subsubsection*{Latihan 1 --- Funktorialitas kernel, kokernel, dan citra}
\hypertarget{o014.mastery.diagram-chasing.001.ex01}{ }
\label{mastery:diagram-chasing-001-ex01}

Misalkan
\[
\begin{tikzcd}
	A \arrow[r, "f"] \arrow[d, "a"'] & B \arrow[d, "b"] \\
	A' \arrow[r, "g"'] & B'
\end{tikzcd}
\qquad (ga=bf)
\]
merupakan persegi komutatif modul-$R$.
\begin{enumerate}[(i)]
	\item Bangun homomorfisme alami
	$\overline a:\Ker(f)\to\Ker(g)$ dan
	$\overline b:\Coker(f)\to\Coker(g)$.
	\item Buktikan $b(\Image(f))\subset\Image(g)$, serta bahwa
	kesamaan berlaku jika $a$ surjektif.
	\item Buktikan $\Ker(f)=a^{-1}(\Ker(g))$ jika $b$ injektif.
\end{enumerate}

% stable-solution-id: o014.mastery.diagram-chasing.001.sol01
\paragraph{Solusi.}
\hypertarget{o014.mastery.diagram-chasing.001.sol01}{ }
Jika $x\in\Ker(f)$, maka $g(a(x))=b(f(x))=0$; jadi
$a(x)\in\Ker(g)$. Pembatasan $a$ memberikan
$\overline a:\Ker(f)\to\Ker(g)$. Untuk kokernel, definisikan
\[
	\overline b([y])=[b(y)], \qquad [y]\in B/\Image(f).
\]
Jika $y$ diganti oleh $y+f(x)$, maka
$b(y+f(x))=b(y)+g(a(x))$, yang menentukan kelas yang sama modulo
$\Image(g)$. Jadi $\overline b$ terdefinisi dengan baik. Kedua konstruksi
ini tepat merupakan konstruksi dalam \eqref{eqn:Ker-Coker-functoriality}.

Selanjutnya,
\[
	b(\Image(f))=b(f(A))=g(a(A))\subset g(A')=\Image(g).
\]
Jika $a$ surjektif, maka $a(A)=A'$, sehingga inklusi tersebut menjadi
kesamaan. Terakhir, $x\in\Ker(f)$ selalu menyiratkan
$a(x)\in\Ker(g)$. Sebaliknya, jika $a(x)\in\Ker(g)$, maka
$0=g(a(x))=b(f(x))$; injektivitas $b$ memberikan $f(x)=0$. Jadi
$x\in\Ker(f)$, dan diperoleh
$\Ker(f)=a^{-1}(\Ker(g))$.

% stable-exercise-id: o014.mastery.diagram-chasing.001.ex02
\subsubsection*{Latihan 2 --- Irisan, jumlah, dan teorema isomorfisme}
\hypertarget{o014.mastery.diagram-chasing.001.ex02}{ }
\label{mastery:diagram-chasing-001-ex02}

Misalkan $U$ dan $V$ submodul dari modul-$R$ $M$. Definisikan
\[
	r:U\cap V\longrightarrow U\oplus V,\quad r(w)=(w,-w),
	\qquad
	s:U\oplus V\longrightarrow U+V,\quad s(u,v)=u+v.
\]
Buktikan bahwa
\[
	0\longrightarrow U\cap V\xrightarrow{r}U\oplus V
	\xrightarrow{s}U+V\longrightarrow0
\]
eksak. Gunakan pengejaran kernel dan citra untuk memperoleh isomorfisme
kanonik
\[
	U/(U\cap V)\rightiso (U+V)/V.
\]

% stable-solution-id: o014.mastery.diagram-chasing.001.sol02
\paragraph{Solusi.}
\hypertarget{o014.mastery.diagram-chasing.001.sol02}{ }
Jika $r(w)=0$, maka komponen pertama menunjukkan $w=0$; jadi $r$ injektif.
Selain itu, $s(r(w))=w-w=0$, sehingga $\Image(r)\subset\Ker(s)$. Jika
$(u,v)\in\Ker(s)$, maka $u+v=0$. Jadi $u=-v$ terletak sekaligus dalam
$U$ dan $V$, serta
\[
	(u,v)=(u,-u)=r(u).
\]
Dengan demikian $\Ker(s)=\Image(r)$. Setiap unsur $U+V$ berbentuk $u+v$,
jadi $s$ surjektif. Barisan tersebut eksak.

Sekarang tinjau homomorfisme
\[
	\theta:U\longrightarrow (U+V)/V,\qquad u\longmapsto u+V.
\]
Pemetaan ini surjektif karena $(u+v)+V=u+V$. Kernelnya terdiri atas semua
$u\in U$ yang juga terletak dalam $V$, yakni $U\cap V$. Faktorisasi
epi--mono, atau Teorema Isomorfisme~\ref{prop:Abel-cat-isom-thm}, memberikan
isomorfisme kanonik
$U/(U\cap V)\rightiso (U+V)/V$.

% stable-exercise-id: o014.mastery.diagram-chasing.001.ex03
\subsubsection*{Latihan 3 --- Lema Lima Pendek}
\hypertarget{o014.mastery.diagram-chasing.001.ex03}{ }
\label{mastery:diagram-chasing-001-ex03}

Dalam diagram komutatif
\[
\begin{tikzcd}
	0 \arrow[r] & A' \arrow[r, "i"] \arrow[d, "a'"'] & A \arrow[r, "p"] \arrow[d, "a"'] & A'' \arrow[r] \arrow[d, "a''"] & 0 \\
	0 \arrow[r] & B' \arrow[r, "j"'] & B \arrow[r, "q"'] & B'' \arrow[r] & 0,
\end{tikzcd}
\]
andaikan kedua baris eksak.
\begin{enumerate}[(i)]
	\item Jika $a'$ dan $a''$ injektif, buktikan bahwa $a$ injektif.
	\item Jika $a'$ dan $a''$ surjektif, buktikan bahwa $a$ surjektif.
	\item Simpulkan bahwa jika $a'$ dan $a''$ isomorfisme, maka $a$ juga
	isomorfisme.
\end{enumerate}

% stable-solution-id: o014.mastery.diagram-chasing.001.sol03
\paragraph{Solusi.}
\hypertarget{o014.mastery.diagram-chasing.001.sol03}{ }
Untuk (i), ambil $x\in A$ dengan $a(x)=0$. Maka
$a''p(x)=qa(x)=0$. Karena $a''$ injektif, $p(x)=0$; keeksakan baris atas
memberi $x=i(x')$. Selanjutnya
$0=a i(x')=j a'(x')$. Karena $j$ dan $a'$ injektif, $x'=0$, sehingga
$x=0$.

Untuk (ii), ambil $y\in B$ dan nyatakan $y''=q(y)$. Surjektivitas $a''$
memberi $x''\in A''$ dengan $a''(x'')=y''$. Karena $p$ surjektif, pilih
$x\in A$ dengan $p(x)=x''$. Sekarang
\[
	q(y-a(x))=y''-a''p(x)=0.
\]
Keeksakan baris bawah memberi $y-a(x)=j(y')$ untuk suatu $y'\in B'$.
Surjektivitas $a'$ memberi $x'\in A'$ dengan $a'(x')=y'$. Maka
\[
	y=a(x)+j(a'(x'))=a(x+i(x')).
\]
Jadi $a$ surjektif. Jika kedua morfisme luar isomorfisme, (i) dan (ii)
menunjukkan bahwa $a$ sekaligus injektif dan surjektif, sehingga merupakan
isomorfisme.

% stable-exercise-id: o014.mastery.diagram-chasing.001.ex04
\subsubsection*{Latihan 4 --- Lema Sembilan}
\hypertarget{o014.mastery.diagram-chasing.001.ex04}{ }
\label{mastery:diagram-chasing-001-ex04}

Perhatikan diagram komutatif modul-$R$
\[
\begin{tikzcd}[column sep=small]
	& 0 \arrow[d] & 0 \arrow[d] & 0 \arrow[d] & \\
	0 \arrow[r] & A' \arrow[r, "f'"] \arrow[d, "\alpha"'] & B' \arrow[r, "g'"] \arrow[d, "\beta"'] & C' \arrow[r] \arrow[d, "\gamma"'] & 0 \\
	0 \arrow[r] & A \arrow[r, "f"] \arrow[d, "\overline\alpha"'] & B \arrow[r, "g"] \arrow[d, "\overline\beta"'] & C \arrow[r] \arrow[d, "\overline\gamma"'] & 0 \\
	0 \arrow[r] & A'' \arrow[r, "f''"] \arrow[d] & B'' \arrow[r, "g''"] \arrow[d] & C'' \arrow[r] \arrow[d] & 0 \\
	& 0 & 0 & 0 &
\end{tikzcd}
\]
dengan ketiga kolom eksak. Andaikan baris atas dan baris tengah eksak.
Buktikan bahwa baris bawah juga eksak.

% stable-solution-id: o014.mastery.diagram-chasing.001.sol04
\paragraph{Solusi.}
\hypertarget{o014.mastery.diagram-chasing.001.sol04}{ }
Pertama kita buktikan bahwa $f''$ injektif. Ambil $a''\in A''$ dengan
$f''(a'')=0$, lalu angkat $a''$ menjadi $a\in A$. Karena
$\overline\beta f(a)=f''\overline\alpha(a)=0$, terdapat $b'\in B'$ dengan
$\beta(b')=f(a)$. Menerapkan $g$ memberi
\[
	\gamma(g'(b'))=g(\beta(b'))=g(f(a))=0.
\]
Morfisme $\gamma$ injektif, sehingga $g'(b')=0$. Keeksakan baris atas
memberi $b'=f'(a')$ untuk suatu $a'\in A'$. Maka
$f(a-\alpha(a'))=0$. Karena $f$ injektif, $a=\alpha(a')$, sehingga
$a''=\overline\alpha(a)=0$.

Berikutnya, $g''$ surjektif. Untuk $c''\in C''$, pilih $c\in C$ yang
dipetakan ke $c''$. Surjektivitas $g$ memberi $b\in B$ dengan $g(b)=c$; maka
$g''(\overline\beta(b))=\overline\gamma(g(b))=c''$.

Tinggal memeriksa keeksakan di $B''$. Komutativitas langsung memberi
$g''f''=0$. Sebaliknya, ambil $b''\in B''$ dengan $g''(b'')=0$ dan angkat
$b''$ ke $b\in B$. Karena
$\overline\gamma(g(b))=g''(\overline\beta(b))=0$, terdapat $c'\in C'$
dengan $\gamma(c')=g(b)$. Pilih $b'\in B'$ dengan $g'(b')=c'$. Maka
\[
	g(b-\beta(b'))=g(b)-\gamma(g'(b'))=0.
\]
Keeksakan baris tengah memberi $a\in A$ dengan
$f(a)=b-\beta(b')$. Setelah dipetakan ke baris bawah, diperoleh
\[
	f''(\overline\alpha(a))
	=\overline\beta(f(a))
	=\overline\beta(b)=b''.
\]
Jadi $\Ker(g'')=\Image(f'')$, dan baris bawah eksak.

% stable-exercise-id: o014.mastery.diagram-chasing.001.ex05
\subsubsection*{Latihan 5 --- Membangun morfisme penghubung}
\hypertarget{o014.mastery.diagram-chasing.001.ex05}{ }
\label{mastery:diagram-chasing-001-ex05}

Kembali ke diagram barisan eksak pendek
\[
\begin{tikzcd}
	0 \arrow[r] & A' \arrow[r, "i"] \arrow[d, "\alpha'"'] & A \arrow[r, "p"] \arrow[d, "\alpha"'] & A'' \arrow[r] \arrow[d, "\alpha''"] & 0 \\
	0 \arrow[r] & B' \arrow[r, "j"'] & B \arrow[r, "q"'] & B'' \arrow[r] & 0.
\end{tikzcd}
\]
Untuk $x''\in\Ker(\alpha'')$, pilih $x\in A$ dengan $p(x)=x''$.
Tunjukkan bahwa $\alpha(x)=j(y')$ untuk suatu $y'\in B'$, lalu definisikan
\[
	\delta(x''):=[y']\in\Coker(\alpha')=B'/\Image(\alpha').
\]
Buktikan bahwa $\delta$ terdefinisi dengan baik, merupakan homomorfisme, dan
alami terhadap morfisme antara dua diagram semacam ini.

% stable-solution-id: o014.mastery.diagram-chasing.001.sol05
\paragraph{Solusi.}
\hypertarget{o014.mastery.diagram-chasing.001.sol05}{ }
Karena $x''\in\Ker(\alpha'')$,
\[
	q(\alpha(x))=\alpha''(p(x))=\alpha''(x'')=0.
\]
Keeksakan baris bawah memberi $\alpha(x)\in\Image(j)$, jadi terdapat
$y'\in B'$ dengan $j(y')=\alpha(x)$. Unsur $y'$ unik karena $j$ injektif.

Jika $x_1$ merupakan pengangkatan lain dari $x''$, maka
$x_1-x\in\Ker(p)=\Image(i)$; tuliskan $x_1-x=i(a')$. Jika
$j(y_1')=\alpha(x_1)$, komutativitas memberi
\[
	j(y_1'-y'-\alpha'(a'))
	=\alpha(x_1)-\alpha(x)-\alpha i(a')=0.
\]
Injektivitas $j$ memberi $y_1'-y'=\alpha'(a')$. Jadi
$[y_1']=[y']$ dalam $\Coker(\alpha')$, dan $\delta$ tidak bergantung pada
pilihan $x$. Untuk $x_1''$ dan $x_2''$, jumlah pengangkatan yang dipilih dan
jumlah unsur $y'$ yang bersesuaian merupakan data yang sah bagi
$x_1''+x_2''$; maka
$\delta$ aditif dan $R$-linear.

Untuk kealamian, ambil morfisme dari diagram di atas ke diagram serupa yang
semua perseginya komutatif. Jika $x$ dan $y'$ dipilih untuk $x''$, citra
$x$ dan $y'$ dalam diagram kedua merupakan pilihan yang sah bagi citra $x''$.
Karena itu, kedua jalur dalam persegi
\[
\begin{tikzcd}
	\Ker(\alpha'') \arrow[r, "\delta"] \arrow[d] & \Coker(\alpha') \arrow[d] \\
	\Ker(\widetilde\alpha'') \arrow[r, "\widetilde\delta"'] & \Coker(\widetilde\alpha')
\end{tikzcd}
\]
memetakan $x''$ ke kelas dari citra $y'$ yang sama. Persegi tersebut
komutatif, sehingga $\delta$ alami. Ini adalah versi berbasis unsur dari sifat
kanonik dalam Catatan~\ref{rem:connecting-canonical}.

% stable-exercise-id: o014.mastery.diagram-chasing.001.ex06
\subsubsection*{Latihan 6 --- Membuktikan Lema Ular}
\hypertarget{o014.mastery.diagram-chasing.001.ex06}{ }
\label{mastery:diagram-chasing-001-ex06}

Dengan diagram dan $\delta$ dari Latihan 5, buktikan bahwa barisan
\[
\begin{aligned}
	0 \longrightarrow \Ker(\alpha') \longrightarrow \Ker(\alpha)
	&\longrightarrow \Ker(\alpha'') \xrightarrow{\delta}
	\Coker(\alpha') \\
	&\longrightarrow \Coker(\alpha)
	\longrightarrow \Coker(\alpha'') \longrightarrow 0
\end{aligned}
\]
eksak. Nyatakan dengan jelas di mana keeksakan kedua baris digunakan dalam
pengejaran.

% stable-solution-id: o014.mastery.diagram-chasing.001.sol06
\paragraph{Solusi.}
\hypertarget{o014.mastery.diagram-chasing.001.sol06}{ }
Semua panah selain $\delta$ diinduksi oleh $i,p,j,q$. Pemetaan
$\Ker(\alpha')\to\Ker(\alpha)$ injektif karena $i$ injektif.

Di $\Ker(\alpha)$, ambil $x\in\Ker(\alpha)$ yang memetakan ke nol dalam
$\Ker(\alpha'')$. Maka $p(x)=0$, sehingga keeksakan baris atas memberi
$x=i(x')$. Persamaan
$0=\alpha(x)=j\alpha'(x')$ dan injektivitas $j$ memberi
$x'\in\Ker(\alpha')$. Inklusi sebaliknya langsung mengikuti dari $pi=0$.
Jadi barisan eksak
di $\Ker(\alpha)$.

Di $\Ker(\alpha'')$, unsur yang berasal dari $x\in\Ker(\alpha)$ mempunyai
$y'=0$ dalam konstruksi Latihan 5, sehingga dipetakan ke nol oleh $\delta$.
Sebaliknya, andaikan $x''\in\Ker(\alpha'')$ dan $\delta(x'')=0$. Pilih
$x\in A$ dan $y'\in B'$ seperti dalam konstruksi $\delta$. Karena
$[y']=0$ dalam $\Coker(\alpha')$, terdapat $a'\in A'$ dengan
$y'=\alpha'(a')$. Maka
\[
	\alpha(x-i(a'))=j(y')-j\alpha'(a')=0,
	\qquad p(x-i(a'))=x''.
\]
Jadi $x''$ berasal dari $\Ker(\alpha)$.

Di $\Coker(\alpha')$, panah berikutnya memetakan $[y']$ ke $[j(y')]$.
Jika $[y']=\delta(x'')$ dan $x$ merupakan pengangkatan yang dipakai, maka
$j(y')=\alpha(x)$, sehingga $[j(y')]=0$ dalam $\Coker(\alpha)$. Sebaliknya,
jika $[j(y')]=0$, terdapat $x\in A$ dengan $j(y')=\alpha(x)$. Komutativitas
memberi
\[
	\alpha''p(x)=q\alpha(x)=qj(y')=0.
\]
Jadi $p(x)\in\Ker(\alpha'')$, dan konstruksi tersebut langsung memberikan
$\delta(p(x))=[y']$.

Di $\Coker(\alpha)$, citra setiap $[y']\in\Coker(\alpha')$ dipetakan ke
$[qj(y')]=0$. Sebaliknya, andaikan $[b]\in\Coker(\alpha)$ dipetakan ke nol.
Artinya, $q(b)=\alpha''(x'')$ untuk suatu $x''\in A''$. Surjektivitas $p$
memberi $x\in A$ dengan $p(x)=x''$. Maka
\[
	q(b-\alpha(x))=q(b)-\alpha''p(x)=0.
\]
Keeksakan baris bawah memberi $b-\alpha(x)=j(y')$ untuk suatu $y'\in B'$.
Karena $\alpha(x)$ lenyap dalam $\Coker(\alpha)$, kelas $[b]$ merupakan
citra $[y']$.

Terakhir, $\Coker(\alpha)\to\Coker(\alpha'')$ surjektif karena $q$ sendiri
surjektif: angkat wakil mana pun di $B''$ ke $B$. Keeksakan di setiap suku
telah diperiksa. Inilah Lema Ular
\ref{prop:snake-lemma} untuk dua barisan eksak pendek modul.

% stable-exercise-id: o014.mastery.diagram-chasing.001.ex07
\subsubsection*{Latihan 7 --- Menghitung ular secara eksplisit}
\hypertarget{o014.mastery.diagram-chasing.001.ex07}{ }
\label{mastery:diagram-chasing-001-ex07}

Terapkan Lema Ular pada diagram grup abelian
\[
\begin{tikzcd}
	0 \arrow[r] & \Z \arrow[r, "\times 6"] \arrow[d, "\times 4"'] & \Z \arrow[r, "\pi"] \arrow[d, "\times 4"'] & \Z/6\Z \arrow[r] \arrow[d, "\times 4"] & 0 \\
	0 \arrow[r] & \Z \arrow[r, "\times 6"'] & \Z \arrow[r, "\pi"'] & \Z/6\Z \arrow[r] & 0.
\end{tikzcd}
\]
Hitung semua kernel dan kokernel, hitung $\delta$ pada generatornya, lalu
tuliskan dan verifikasi barisan eksak yang dihasilkan.

% stable-solution-id: o014.mastery.diagram-chasing.001.sol07
\paragraph{Solusi.}
\hypertarget{o014.mastery.diagram-chasing.001.sol07}{ }
Kedua pemetaan vertikal pertama adalah perkalian dengan $4$ pada $\Z$, jadi
kernelnya nol dan kokernelnya $\Z/4\Z$. Pada kolom ketiga,
\[
	\Ker(\times4:\Z/6\Z\to\Z/6\Z)=\{[0],[3]\}\simeq\Z/2\Z.
\]
Citra perkalian dengan $4$ pada $\Z/6\Z$ ialah
$\{[0],[2],[4]\}$, sehingga kokernelnya isomorfik dengan $\Z/2\Z$.

Untuk menghitung $\delta([3])$, angkat $[3]$ menjadi $3\in\Z$. Pemetaan
vertikal tengah mengirimkannya ke $12$. Karena $12=6\cdot2$, konstruksi
morfisme penghubung memberi
\[
	\delta([3])=[2]\in\Z/4\Z.
\]
Dengan identifikasi generator $[3]\leftrightarrow[1]\in\Z/2\Z$, pemetaan
$\delta$ ialah $[1]\mapsto[2]$.

Pemetaan dari kokernel pertama ke kokernel kedua diinduksi oleh
$\times6$, sehingga pada $\Z/4\Z$ sama dengan $\times2$. Pemetaan berikutnya
ialah reduksi modulo $2$. Jadi barisan Ular menjadi
\[
	0\longrightarrow0\longrightarrow0\longrightarrow\Z/2\Z
	\mathop{\longrightarrow}^{[1]\mapsto[2]}\Z/4\Z
	\xrightarrow{\times2}\Z/4\Z
	\xrightarrow{\bmod 2}\Z/2\Z\longrightarrow0.
\]
Citra $\delta$ ialah $\{[0],[2]\}=\Ker(\times2)$, sedangkan citra
$\times2$ juga $\{[0],[2]\}=\Ker(\bmod 2)$. Pemetaan terakhir surjektif,
dan $\delta$ injektif. Jadi barisan tersebut benar-benar eksak.

% stable-exercise-id: o014.mastery.diagram-chasing.001.ex08
\subsubsection*{Latihan 8 --- Morfisme penghubung dalam kohomologi}
\hypertarget{o014.mastery.diagram-chasing.001.ex08}{ }
\label{mastery:diagram-chasing-001-ex08}

Misalkan
\[
	0\longrightarrow A^\bullet\xrightarrow{i}B^\bullet
	\xrightarrow{q}C^\bullet\longrightarrow0
\]
merupakan barisan eksak pendek kompleks kokrantai modul-$R$, eksak pada
setiap derajat. Untuk kelas $[c]\in\Hm^n(C^\bullet)$, pilih wakil kokitar
$c\in C^n$ dan pengangkatan $b\in B^n$ dengan $q(b)=c$.
\begin{enumerate}[(i)]
	\item Buktikan bahwa terdapat $a\in A^{n+1}$ dengan
	$i(a)=d_B(b)$, dan bahwa $a$ merupakan kokitar.
	\item Buktikan bahwa rumus $\partial^n[c]=[a]$ memberikan homomorfisme
	yang terdefinisi dengan baik
	$\partial^n:\Hm^n(C^\bullet)\to\Hm^{n+1}(A^\bullet)$.
	\item Buktikan keeksakan ruas
	\[
		\Hm^n(A^\bullet)\longrightarrow\Hm^n(B^\bullet)
		\longrightarrow\Hm^n(C^\bullet)\xrightarrow{\partial^n}
		\Hm^{n+1}(A^\bullet)\longrightarrow\Hm^{n+1}(B^\bullet).
	\]
	\item Buktikan bahwa $\partial^n$ alami terhadap morfisme barisan eksak
	pendek kompleks.
\end{enumerate}

% stable-solution-id: o014.mastery.diagram-chasing.001.sol08
\paragraph{Solusi.}
\hypertarget{o014.mastery.diagram-chasing.001.sol08}{ }
Karena $c$ kokitar,
\[
	q(d_Bb)=d_C(qb)=d_Cc=0.
\]
Keeksakan pada $B^{n+1}$ memberi
$d_Bb\in\Image(i^{n+1})$, sehingga ada $a\in A^{n+1}$ dengan
$i(a)=d_Bb$. Selanjutnya,
\[
	i(d_Aa)=d_B(i(a))=d_B^2b=0.
\]
Karena $i$ injektif pada derajat $n+2$, diperoleh $d_Aa=0$; jadi $a$
kokitar.

Jika pengangkatan $b$ diganti oleh $b+i(u)$, maka unsur yang bersesuaian
berubah dari $a$ menjadi $a+d_Au$, sehingga kelas $[a]$ tidak berubah.
Jika wakil $c$ diganti oleh $c+d_Cv$, pilih pengangkatan $w\in B^{n-1}$
dari $v$ dan gunakan $b+d_Bw$ sebagai pengangkatan wakil baru. Karena
$d_B(b+d_Bw)=d_Bb$, kelas hasilnya tetap $[a]$. Khususnya, jika $c$ suatu
kokobatas, kita dapat memilih pengangkatan berbentuk $d_Bw$ dan memperoleh
$a=0$. Jadi $\partial^n$ bergantung hanya pada kelas kohomologi. Penggunaan
jumlah pengangkatan menunjukkan bahwa $\partial^n$ merupakan homomorfisme.

Kita periksa keeksakan pada setiap suku. Di $\Hm^n(B^\bullet)$,
citra kelas yang berasal dari $A$ jelas lenyap setelah diterapkan $q$.
Sebaliknya, jika
$b\in B^n$ kokitar dan $[q(b)]=0$, tuliskan
$q(b)=d_Cv$ dan angkat $v$ menjadi $w\in B^{n-1}$. Maka
$b-d_Bw\in\Ker(q)=\Image(i)$; tuliskan $b-d_Bw=i(a)$. Karena $b-d_Bw$
kokitar dan $i$ injektif, $a$ juga kokitar. Jadi $[b]$ berasal dari
$\Hm^n(A^\bullet)$.

Di $\Hm^n(C^\bullet)$, kelas citra dari kokitar $b\in B^n$ mempunyai
$d_Bb=0$, sehingga dipetakan ke nol oleh $\partial^n$. Sebaliknya, jika
$\partial^n[c]=0$, konstruksi memberi $i(a)=d_Bb$ dengan $a=d_Au$ untuk
suatu $u\in A^n$. Maka
\[
	d_B(b-i(u))=i(a)-i(d_Au)=0,
	\qquad q(b-i(u))=c.
\]
Jadi $[c]$ berasal dari $\Hm^n(B^\bullet)$.

Di $\Hm^{n+1}(A^\bullet)$, kelas $\partial^n[c]=[a]$ dipetakan ke
$[i(a)]=[d_Bb]=0$. Sebaliknya, andaikan $a\in A^{n+1}$ kokitar dan
$[i(a)]=0$ dalam $\Hm^{n+1}(B^\bullet)$. Pilih $b\in B^n$ dengan
$i(a)=d_Bb$. Unsur $c=q(b)$ merupakan kokitar karena
$d_Cc=q(d_Bb)=q(i(a))=0$, dan konstruksi memberi
$\partial^n[c]=[a]$. Bagian barisan yang diminta pun eksak.

Terakhir, ambil morfisme antara dua barisan eksak pendek kompleks. Jika
$c,b,a$ merupakan data pengejaran di atas, citra ketiganya di barisan kedua
masih memenuhi
\[
	q(b)=c,\qquad i(a)=d_Bb.
\]
Karena itu, menerapkan morfisme lebih dahulu lalu membentuk morfisme
penghubung memberikan kelas yang sama dengan membentuk $\partial^n$ lebih
dahulu lalu menerapkan morfisme pada kohomologi. Jadi $\partial^n$ alami.

Dengan menjalankan konstruksi ini pada setiap $n$, ruas-ruas eksak tersebut
menyambung menjadi barisan eksak panjang
\[
	\cdots\to\Hm^n(A^\bullet)\to\Hm^n(B^\bullet)\to\Hm^n(C^\bullet)
	\xrightarrow{\partial^n}\Hm^{n+1}(A^\bullet)\to\Hm^{n+1}(B^\bullet)
	\to\cdots.
\]
Morfisme $\partial^n$ adalah morfisme penghubung yang sama secara konseptual
dengan morfisme pada Lema Ular: sebuah kokitar diangkat, diferensialnya
didorong, lalu hasil yang kini berada dalam kernel diturunkan ke suku kiri.
